% =========================================================================
% Capítulo 1: Cinemática
% Tema 05: Problemas de mezcla compleja
% Problema 03: Mecanismo de rotación tridimensional
% =========================================================================
\subsubsection{Mecanismo de rotación tridimensional}
Un sistema mecánico consiste en un eje vertical que gira con una velocidad angular $\omega_1$. Unido a este eje en el punto $O$, se encuentra un brazo de longitud $L$ que forma un ángulo $\beta$ con la vertical y gira sobre su propio eje con una velocidad angular $\omega_2$. En el extremo del brazo se sitúa un disco de radio $r$ que rota con una velocidad angular $\omega_3$ respecto al brazo, como se ilustra en la Figura \ref{fig:mecanismo_3d}. Utilizando el sistema de referencia indicado, determine:

\begin{enumerate}
    \item El vector velocidad angular total del disco respecto al suelo.
    \item La velocidad y aceleración del punto $P$ situado en el borde del disco.
    \item La aceleración del centro del disco $C$ y del punto inferior $D$.
\end{enumerate}

\begin{figure}[!ht]
    \centering
    \begin{tikzpicture}[>=Latex, scale=1.5]
        % --- Configuración de Ejes y Referencias ---
        \coordinate (O) at (0,0);
        
        % Línea de eje vertical (punteada)
        \draw[dotted, thick] (0,1) -- (0,-4);
        
        % Ejes e_x y e_z (azules)


        % --- Soporte Superior ---
        \shadedraw[top color=gray!60, bottom color=gray!10] (-0.2,0) rectangle (0.2,1.2);
        \draw[thick] (-0.2,0.5) -- (0.2,0.5);
        
        % Velocidad angular omega_1
        \draw[->, thick] (0.3,0.9) arc (0:150:0.3 and 0.1);
        \node at (0.5,1) {$\omega_1$};
        
        % Punto O
        \filldraw (O) circle (2pt) node[above left] {$O$};

        % --- El Brazo (Rod) ---
        % Ángulo beta y longitud L
        \def\betaAng{35} % Ángulo respecto a la vertical
        \def\L{4.2}      % Longitud del brazo
        
        % Coordenada final del brazo (Punto C)
        % Nota: 270 es abajo, sumamos betaAng para mover a la derecha
        \coordinate (C) at ({\betaAng-90}:\L);
        
        % Dibujo del brazo cilíndrico
        \begin{scope}[rotate around={\betaAng:(O)}]
            \shadedraw[left color=gray!50, right color=gray!20, middle color=gray!5,line cap=round ] 
                (-0.1,0) rectangle (0.1,-\L);
            
            % Velocidad angular omega_2
            \draw[<-, thick] (0.3, -\L*0.5) arc (0:180:0.3 and 0.1);
            \node at (0.6, -\L*0.5 + 0.2) {$\omega_2$};
        \end{scope}

        % Indicación de longitud L
        \draw[<->] ($(O)+({180+\betaAng}:0.4)$) -- ($(C)+({180+\betaAng}:0.4)$);
        \node[fill=white, inner sep=1pt] at ($(O)!0.5!(C) + ({180+\betaAng}:0.6)$) {$L$};

        % Ángulo Beta
        \draw[<->, thick] (0,-1.2) arc (270:{270+\betaAng}:1.2);
        \node at (0.25,-1) {$\beta$};

        % --- El Disco Terminal ---
        \begin{scope}[shift=(C), rotate around={\betaAng:(0,0)}]
            % Perfil del disco (rectángulo sombreado)
            \shadedraw[top color=gray!10, bottom color=gray!40] (-0.6,-0.2) rectangle (0.6,0.2);
            \draw[thick] (-0.6,-0.2) rectangle (0.6,0.2);
            
            % Puntos del disco
            \filldraw (0,0) circle (1.2pt) node[below left] {$C$}; % Centro
            \filldraw (0.6,0) circle (2pt) node[below right] {$P$}; % Extremo P
            \node[below] at (0,-0.2) {$\mathcal{D}$}; % Punto D
            
            % Radio r
            \draw[|-|] (0,0.4) -- (0.6,0.4) node[midway, above] {$r$};
            \draw[dotted] (0.6,0.2) -- (0.9,0.2);
            \draw[dotted] (0,0) -- (0.9,0);

            % Velocidad angular omega_3
            \draw[->, thick] (-0.4, -0.4) arc (-180:0
            :0.6 and 0.2);
            \node at (0.23,-0.5) {$\omega_3$};
        \end{scope}

    \end{tikzpicture}
    \caption{Esquema de un mecanismo con tres rotaciones superpuestas: rotación del soporte ($\omega_1$), rotación axial del brazo ($\omega_2$) y rotación del disco terminal ($\omega_3$).}
    \label{fig:mecanismo_3d}
\end{figure}